% Bagian: turing-machines
% Bab: machines-computations
% Subbagian: variants

\documentclass[../../../include/open-logic-section]{subfiles}

\begin{document}

\olfileid[id]{tur}{mac}{var}
\olsection{Varian Mesin Turing}

Sesungguhnya ada banyak cara yang mungkin untuk mendefinisikan mesin Turing;
definisi kita hanyalah salah satunya. Dalam beberapa hal, definisi kita lebih
longgar daripada definisi lain. Kita mengizinkan alfabet berhingga sebarang,
sedangkan definisi yang lebih terbatas mungkin hanya mengizinkan dua simbol
pita, $\TMstroke$ dan~$\TMblank$. Kita mengizinkan mesin menuliskan simbol pada
pita sekaligus bergerak, sedangkan definisi lain hanya mengizinkan penulisan
atau pergerakan pada satu waktu. Kita mengizinkan penulisan tanpa menggerakkan
kepala pita, sedangkan definisi lain tidak memiliki ``instruksi''~$\TMstay$.
Dalam hal lain, definisi kita justru lebih membatasi. Kita mengandaikan bahwa
pita hanya tak berhingga dalam satu arah, sedangkan definisi lain mengizinkan
pita tak berhingga ke kiri maupun ke kanan. Bahkan, orang dapat mengizinkan
sebarang banyak pita terpisah, atau kisi petak yang tak berhingga. Kita
merepresentasikan himpunan instruksi mesin Turing dengan fungsi transisi;
definisi lain menggunakan relasi transisi sehingga mesin memiliki lebih dari
satu kemungkinan instruksi dalam suatu keadaan tertentu.

Pelonggaran definisi yang terakhir ini sangat menarik. Dalam definisi kita,
ketika mesin berada dalam keadaan~$q$ dan membaca simbol~$\sigma$,
$\delta(q, \sigma)$ menentukan simbol baru, keadaan baru, dan posisi kepala
pita yang baru. Namun, jika kita mengizinkan himpunan instruksi berupa relasi
antara pasangan keadaan-simbol saat ini $\tuple{q, \sigma}$ dan tripel
keadaan-simbol-arah yang baru $\tuple{q', \sigma', D}$, tindakan mesin Turing
mungkin tidak ditentukan secara tunggal---relasi instruksi itu mungkin memuat
baik $\tuple{q, \sigma, q', \sigma', D}$ maupun
$\tuple{q, \sigma, q'', \sigma'', D'}$. Dalam kasus ini kita memperoleh mesin
Turing \emph{nondeterministik}. Mesin seperti ini berperan penting dalam teori
kompleksitas komputasi.

Ada pula konvensi yang berbeda mengenai kapan mesin Turing berhenti. Menurut
kita, mesin berhenti ketika fungsi transisinya tidak terdefinisi; definisi lain
mengharuskan mesin berada dalam keadaan penghentian khusus. Dalam
\olref[hal]{sec} telah kita jelaskan mengapa persyaratan adanya keadaan
penghentian khusus bukan pembatasan yang memengaruhi apa yang dapat dihitung
oleh mesin Turing. Karena pita mesin Turing kita hanya tak berhingga dalam
satu arah, terdapat kasus ketika mesin Turing tidak dapat menjalankan
instruksi dengan semestinya: ketika mesin membaca petak paling kiri dan harus
bergerak ke kiri. Menurut definisi kita, mesin hanya tetap di tempat alih-alih
``jatuh dari pita'', tetapi kita juga dapat mendefinisikan bahwa mesin berhenti
ketika hal itu terjadi. Definisi ini pun ekuivalen: perilaku mesin Turing yang
berhenti ketika mencoba bergerak ke kiri dari petak~$0$ dapat kita simulasikan
dengan menghapus setiap transisi $\delta(q,\TMendtape) =
\tuple{q',\sigma,\TMleft}$---dengan demikian, alih-alih mencoba bergerak ke
kiri ketika membaca~$\TMendtape$, mesin berhenti.\footnote{Cara ini tidak
\emph{sepenuhnya} berhasil karena tidak ada yang mencegah kita menulis dan
membaca $\TMendtape$ pada petak selain petak~$0$ (lihat
\olref[una]{ex:mover}). Kita dapat mengatasinya dengan menambahkan simbol
kedua~$\TMendtape'$ untuk digunakan bagi keperluan tersebut.}

Ada pula cara yang berbeda untuk merepresentasikan bilangan (dan dengan
demikian fungsi masukan-keluaran yang dihitung mesin Turing). Kita menggunakan
representasi uner, tetapi representasi biner juga dapat digunakan. Hal ini
memerlukan dua simbol selain $\TMblank$ dan~$\TMendtape$.

Sekarang perhatikan fakta menarik berikut: tidak satu pun dari variasi ini
memengaruhi fungsi mana yang dapat dihitung oleh mesin Turing. \emph{Jika suatu
fungsi dapat dihitung oleh mesin Turing menurut salah satu definisi, fungsi itu
dapat dihitung oleh mesin Turing menurut semuanya.}

Kita tidak akan membahas perincian pembuktiannya. Berikut hanya satu contoh:
pita yang tak berhingga dalam kedua arah atau penggunaan beberapa pita tidak
memberikan daya komputasi tambahan. Secara garis besar, alasannya ialah bahwa
mesin Turing dengan satu pita yang tak berhingga ke satu arah dapat
menyimulasikan beberapa pita atau pita yang tak berhingga ke dua arah.
Misalnya, dengan keadaan dan instruksi tambahan, kita dapat ``menerjemahkan''
program bagi mesin dengan beberapa pita atau pita yang tak berhingga ke dua
arah menjadi program bagi mesin dengan satu pita yang tak berhingga ke satu
arah. Mesin hasil terjemahan dapat menggunakan petak genap untuk petak-petak
pita~$1$ (atau petak ``positif'' pada pita yang tak berhingga ke dua arah) dan
petak ganjil untuk petak-petak pita~$2$ (atau petak ``negatif'').

\end{document}
